\chapter{10 truyện cực ngắn ám ảnh}
\markboth{10 truyện cực ngắn ám ảnh}{10 Haunting Short Stories}


\section{Chiếc ghế trống}
\begin{truyen}{Chiếc ghế trống}{Life lesson}
\chuhoa{M}{ỗi năm ngày giỗ, ông vẫn dọn thêm một chỗ. Chiếc ghế trống. Con cháu hỏi, ông không nói. Đó là ghế của đứa con ông mất từ lâu.}
\textit{(Every year on the death anniversary he still set an extra place. The empty chair. His grandchildren asked; he didn't answer. It was for the child he had lost long ago.)}
\end{truyen}

\begin{baihoc}
Mất con là vết thương không bao giờ lành.
\textit{(Losing a child is a wound that never heals.)}
\end{baihoc}

\begin{ghinhoanh}
\textbf{Short English to remember:}

Losing a child is a wound that never heals.

\end{ghinhoanh}

\begin{tuhocnhanh}
1. Câu chuyện này gợi cho bạn suy nghĩ gì? \textit{What does this story make you think?}

2. Bạn sẽ nhớ bài học nào nhất? \textit{Which lesson will you remember most?}
\end{tuhocnhanh}
\ngancach


\section{Tin nhắn chưa gửi}
\begin{truyen}{Tin nhắn chưa gửi}{Life lesson}
\chuhoa{Đ}{iện thoại bà còn lưu tin nhắn soạn cho con: ``Mẹ nhớ con.'' Chưa kịp gửi thì bà qua đời. Con bà đọc được khi đã muộn.}
\textit{(The phone still had the draft message to her child: ``Mom misses you.'' She hadn't sent it when she passed. Her child read it when it was too late.)}
\end{truyen}

\begin{baihoc}
Đừng để lời yêu thương đến muộn.
\textit{(Don't let words of love come too late.)}
\end{baihoc}

\begin{ghinhoanh}
\textbf{Short English to remember:}

Don't let words of love come too late.

\end{ghinhoanh}

\begin{tuhocnhanh}
1. Câu chuyện này gợi cho bạn suy nghĩ gì? \textit{What does this story make you think?}

2. Bạn sẽ nhớ bài học nào nhất? \textit{Which lesson will you remember most?}
\end{tuhocnhanh}
\ngancach


\section{Hai bát cơm}
\begin{truyen}{Hai bát cơm}{Life lesson}
\chuhoa{Ô}{ng cứ đặt hai bát mỗi bữa. Một cho mình, một cho vợ đã mất. Ông ăn một mình bên hai bát cơm.}
\textit{(He always set two bowls at each meal. One for him, one for his late wife. He ate alone beside two bowls of rice.)}
\end{truyen}

\begin{baihoc}
Tình nghĩa vợ chồng sống mãi.
\textit{(Marriage love lives on.)}
\end{baihoc}

\begin{ghinhoanh}
\textbf{Short English to remember:}

Marriage love lives on.

\end{ghinhoanh}

\begin{tuhocnhanh}
1. Câu chuyện này gợi cho bạn suy nghĩ gì? \textit{What does this story make you think?}

2. Bạn sẽ nhớ bài học nào nhất? \textit{Which lesson will you remember most?}
\end{tuhocnhanh}
\ngancach


\section{Cánh cửa không đóng}
\begin{truyen}{Cánh cửa không đóng}{Life lesson}
\chuhoa{Đ}{êm nào anh cũng để hé cửa. Vì ngày xưa con anh đi hoang, một đêm về thì cửa đã khoá. Từ đó anh không bao giờ đóng.}
\textit{(Every night he left the door ajar. Because long ago his child ran away, and one night came back to a locked door. Since then he never closed it.)}
\end{truyen}

\begin{baihoc}
Đừng đóng cửa với người muốn về.
\textit{(Don't close the door on those who want to return.)}
\end{baihoc}

\begin{ghinhoanh}
\textbf{Short English to remember:}

Don't close the door on those who want to return.

\end{ghinhoanh}

\begin{tuhocnhanh}
1. Câu chuyện này gợi cho bạn suy nghĩ gì? \textit{What does this story make you think?}

2. Bạn sẽ nhớ bài học nào nhất? \textit{Which lesson will you remember most?}
\end{tuhocnhanh}
\ngancach


\section{Cuốn sổ nợ}
\begin{truyen}{Cuốn sổ nợ}{Life lesson}
\chuhoa{B}{à giữ cuốn sổ ghi tên từng người bà đã giúp. Không phải để đòi — để khi bà mất, con cháu biết đi trả ơn.}
\textit{(She kept a notebook of everyone she had helped. Not to collect — so when she died, her children would know to repay the kindness.)}
\end{truyen}

\begin{baihoc}
Ơn nghĩa cần trả.
\textit{(Kindness deserves to be repaid.)}
\end{baihoc}

\begin{ghinhoanh}
\textbf{Short English to remember:}

Kindness deserves to be repaid.

\end{ghinhoanh}

\begin{tuhocnhanh}
1. Câu chuyện này gợi cho bạn suy nghĩ gì? \textit{What does this story make you think?}

2. Bạn sẽ nhớ bài học nào nhất? \textit{Which lesson will you remember most?}
\end{tuhocnhanh}
\ngancach


\section{Tiếng gọi trong đêm}
\begin{truyen}{Tiếng gọi trong đêm}{Life lesson}
\chuhoa{Ô}{ng già hay gọi tên vợ lúc nửa đêm. Tỉnh dậy mới biết bà đã mất mười năm. Thói quen vẫn còn.}
\textit{(The old man often called his wife's name at midnight. He woke to remember she had been gone ten years. The habit remained.)}
\end{truyen}

\begin{baihoc}
Thói quen yêu thương không chết.
\textit{(The habit of love doesn't die.)}
\end{baihoc}

\begin{ghinhoanh}
\textbf{Short English to remember:}

The habit of love doesn't die.

\end{ghinhoanh}

\begin{tuhocnhanh}
1. Câu chuyện này gợi cho bạn suy nghĩ gì? \textit{What does this story make you think?}

2. Bạn sẽ nhớ bài học nào nhất? \textit{Which lesson will you remember most?}
\end{tuhocnhanh}
\ngancach


\section{Bức thư chưa mở}
\begin{truyen}{Bức thư chưa mở}{Life lesson}
\chuhoa{S}{au khi cha mất, con tìm thấy bức thư đề ``Gửi con''. Cha viết từ năm trước. Con không dám mở — sợ nghe giọng cha lần cuối.}
\textit{(After his father died he found a letter addressed ``To my child''. His father had written it the year before. He didn't dare open it — afraid to hear his father's voice one last time.)}
\end{truyen}

\begin{baihoc}
Lời người đã khuất còn đó.
\textit{(The words of the departed remain.)}
\end{baihoc}

\begin{ghinhoanh}
\textbf{Short English to remember:}

The words of the departed remain.

\end{ghinhoanh}

\begin{tuhocnhanh}
1. Câu chuyện này gợi cho bạn suy nghĩ gì? \textit{What does this story make you think?}

2. Bạn sẽ nhớ bài học nào nhất? \textit{Which lesson will you remember most?}
\end{tuhocnhanh}
\ngancach


\section{Chiếc áo cũ}
\begin{truyen}{Chiếc áo cũ}{Life lesson}
\chuhoa{B}{à không nỡ vứt áo của chồng. Bà mặc khi ở nhà. Con bảo bỏ đi. Bà nói: ``Mặc vào như còn ông ấy bên cạnh.''}
\textit{(She couldn't bear to throw away her husband's shirt. She wore it at home. Her children said to get rid of it. She said: ``Wearing it is like having him beside me.'')}
\end{truyen}

\begin{baihoc}
Kỷ vật giữ người ta sống trong lòng.
\textit{(Keepsakes keep them alive in the heart.)}
\end{baihoc}

\begin{ghinhoanh}
\textbf{Short English to remember:}

Keepsakes keep them alive in the heart.

\end{ghinhoanh}

\begin{tuhocnhanh}
1. Câu chuyện này gợi cho bạn suy nghĩ gì? \textit{What does this story make you think?}

2. Bạn sẽ nhớ bài học nào nhất? \textit{Which lesson will you remember most?}
\end{tuhocnhanh}
\ngancach


\section{Bữa cơm cuối}
\begin{truyen}{Bữa cơm cuối}{Life lesson}
\chuhoa{C}{ả nhà bận, năm Tết không về. Ông bà ăn Tết hai người. Năm sau ông mất. Bữa cơm cuối cùng cùng nhau là bữa ấy.}
\textit{(The family was busy, didn't come home for Tet. The old couple had Tet alone. The next year he passed. Their last meal together was that one.)}
\end{truyen}

\begin{baihoc}
Về khi còn có thể.
\textit{(Return while you still can.)}
\end{baihoc}

\begin{ghinhoanh}
\textbf{Short English to remember:}

Return while you still can.

\end{ghinhoanh}

\begin{tuhocnhanh}
1. Câu chuyện này gợi cho bạn suy nghĩ gì? \textit{What does this story make you think?}

2. Bạn sẽ nhớ bài học nào nhất? \textit{Which lesson will you remember most?}
\end{tuhocnhanh}
\ngancach


\section{Ngọn nến}
\begin{truyen}{Ngọn nến}{Life lesson}
\chuhoa{Đ}{êm nào bà cũng thắp một ngọn nến bên ảnh chồng. Bà nói: ``Để ông ấy thấy đường mà về.'' Bà tin linh hồn còn đó.}
\textit{(Every night she lit a candle by her husband's photo. She said: ``So he can see the way home.'' She believed his spirit was still there.)}
\end{truyen}

\begin{baihoc}
Tình yêu vượt qua cái chết.
\textit{(Love transcends death.)}
\end{baihoc}

\begin{ghinhoanh}
\textbf{Short English to remember:}

Love transcends death.

\end{ghinhoanh}

\begin{tuhocnhanh}
1. Câu chuyện này gợi cho bạn suy nghĩ gì? \textit{What does this story make you think?}

2. Bạn sẽ nhớ bài học nào nhất? \textit{Which lesson will you remember most?}
\end{tuhocnhanh}
\ngancach
